\problem{}
% You are provided with three sequences $A$, $B$, and $C$. The lengths of these sequences are $m$, $n$, and $m+n$, respectively. In other words, sequence $C$ has a length that equals the combined lengths of sequences $A$ and $B$. Develop an algorithm to determine if the sequences $A$ and $B$ can be interleaved to form sequence $C$, while maintaining the relative order of characters in $A$ and $B$.

% \textbf{Example 1:} Given $A = aabb$, $B = cba$, and $C = acabbab$, the algorithm should return \textbf{true}.

% \textbf{Example 2:} Given $A = aabb$, $B = cba$, and $C = aaabbbc$, the algorithm should return \textbf{false}.

Determine the time complexity of the pseudocode below and explain your reasoning.

\begin{lstlisting}
def count_triples(arr, threshold):
    n = len(arr)
    count = 0
    
    def quicksort(arr, left, right):
        if left >= right:
            return
        pivot = arr[right]
        i = left
        for j in range(left, right):
            if arr[j] <= pivot:
                arr[i], arr[j] = arr[j], arr[i]
                i += 1
        arr[i], arr[right] = arr[right], arr[i]
        quicksort(arr, left, i - 1)
        quicksort(arr, i + 1, right)
    quicksort(arr, 0, n - 1)
    
    def upper_bound(arr, left, right, target):
        result = left - 1
        while left <= right:
            mid = (left + right) // 2
            if arr[mid] <= target:
                result = mid
                left = mid + 1
            else:
                right = mid - 1
        return result

    for i in range(n):
        for j in range(i + 1, n):
            target = threshold - (arr[i] + arr[j])
            k = upper_bound(arr, j + 1, n - 1, target)
            if k > j:
                count += (k - j)

    return count
\end{lstlisting}

\noindent \textbf{Solution:} \\

\noindent The pseudocode structure can be divided into two parts:
\[
\left\{
\begin{array}{l}
\text{using quicksort on an array of length $n$}  \\
\text{using a double loop that calls \texttt{upper\_bound} inside to calculate k. 
}
\end{array}
\right .
\]


\noindent Thus, the time complexity of the pseudocode =

complexity of quicksort + complexity of calculating k.
\\ 

\noindent For quicksort:

The average-case complexity is $O(n\log n)$, 
where the recursion depth is $O(\log n)$ and the overhead of partition traversal is $O(n)$.

The worst-case complexity is $O(n^2)$, 
in which case the recursion depth is $O(n)$ while the overhead of partition traversal remains $O(n)$.
\\

\noindent For k:

The outer loop is executes n times, the inner loop runs $n-i$ times, and the loop body (calling \texttt{upper\_bound}) has complexity $O(\log n)$.

$ \sum_{i=  1}^{n}\sum_{j=  i+1}^{n}\log_2 n=  \frac{1}{2}n^2 \log_2 n- \frac{1}{2}n \log_2 n$, which is $O(n^2 \log n)$.
\\

\noindent Whether in the worst case or average case, the double loop dominates the overall complexity, and the quicksort term can be omitted.
\\

\noindent Thus, the time complexity of the overall pseudocode is $O(n^2 \log n)$.
\newpage